\section{Introduction}

Most models of lunar settlement are written to answer a question that has
already been posed. A structural analysis asks whether a given span will
stand \citep{blair2017stability}; a thermal model asks what a given shell
costs to keep warm. Such models are precise about a narrow thing.

This work takes the opposite approach. \emph{Lunar Metropolis} is an
executable model of a settlement as a whole --- terrain, illumination,
volatiles, three utility networks, an economy, and a population that responds
to all of them --- built as a playable simulation and then interrogated as a
model. It is deliberately imprecise about every individual mechanism, and in
exchange it can be run end to end, tens of thousands of simulated days at a
time, across hundreds of independently generated worlds. What such a model is
good for is not predicting a number but exposing a \emph{structure}: which
constraints bind, which trade against which, and which supposedly free choices
turn out to have been made by the terrain generator before the player arrived.

The specific question we put to it here is the balance between building on the
lunar surface and building underneath it. The case for going underground is
well rehearsed. Burial under regolith reduces the galactic cosmic ray dose,
though not as simply as one might expect: secondary neutrons, pions and photons
produced within the shield itself first \emph{rise} with shell thickness before
falling away \citep{jia2010radiation}. An intact lava tube offers a large
shielded volume that nobody has to excavate, and the protection is substantial
--- a Monte Carlo study puts the effective dose equivalent inside a horizontal
lunar tube below $1$~mSv\,yr$^{-1}$, against $416$~mSv\,yr$^{-1}$ on the open
surface \citep{naito2020radiation} --- in voids whose morphology at Marius Hills
was mapped from imagery decades ago \citep{greeley1971mariushills} and whose
presence there is supported by radar sounding \citep{kaku2017intact}. What is less often stated is
what the choice \emph{costs}, in a currency that lets it be compared against
simply building on the surface instead. That is the question this model is
shaped to answer.

Two features of the Moon do most of the work in the model, and both are
consequences of illumination. Near the poles the Sun tracks close to the
horizon, so modest relief casts enormous shadows, and crater floors that never
see the Sun act as cold traps where water ice is stable
\citep{paige2010diviner}; the LCROSS impact into Cabeus returned direct
evidence of water in such a deposit \citep{colaprete2010water}. The same
illumination geometry that puts the ice there also makes that ground worthless
to build on: it is dark, cold, and generates no solar power. The model derives
both facts from one raycast, and the tension between them is the engine of
everything reported below.

We make three contributions. First, we describe the model and the reasoning
behind its structure (Sections~\ref{sec:model}--\ref{sec:habitats}). Second, we
calibrate its resource parameters against its own generated worlds rather than
by hand, and report the calibration
(Section~\ref{sec:calibration}). Third, we report what the model says about the
surface/subsurface balance, and what it says about the automated director that
plays it (Sections~\ref{sec:results}--\ref{sec:failures}). We are explicit
throughout about which statements are results of the model and which are
assumptions built into it; Section~\ref{sec:limitations} is devoted to the
difference.
